\documentclass[11pt,a4paper]{article}
\usepackage[margin=1in]{geometry}
\usepackage[hidelinks]{hyperref}
\usepackage{enumitem}
\usepackage{graphicx}
\usepackage{xcolor}
\usepackage{booktabs}
\usepackage{array}
\usepackage{amssymb}  % provides \checkmark, used in the User Roles table

% -----------------------------------------------------
% bvraghav-tiet-ucs503p-article.sty
% -----------------------------------------------------

% Variable: Lab Instructor
\makeatletter
\newcommand{\BvrSetLabInstructor}[1]{\gdef\Bvr@LabInstructor{#1}}
% default (placeholder -- fill in before submission)
\newcommand{\Bvr@LabInstructor}{[INSTRUCTOR NAME]}
% public accessor
\newcommand{\BvrLabInstructorValue}{\Bvr@LabInstructor}
\makeatother

% Add subtitle
\usepackage{titling}
% -- Set title bold
\pretitle{\begin{center}\bfseries\fontsize{18bp}{18bp}\selectfont}
\makeatletter
\newcommand{\BvrSubtitle}[1]{%
  \posttitle{%
    \par\end{center}
    \vspace{-0.6em}
    \begin{center}\large#1\end{center}
    \vskip 0.5em}%
}
\makeatother

% Hints in the section title(s)
\newcommand{\BvrHint}[1]{{\normalfont\normalsize\itshape\hfill #1}}

% -----------------------------------------------------

\title{Civic Issue Reporter}

\BvrSetLabInstructor{[INSTRUCTOR NAME]}

\BvrSubtitle{%
  A Campus Infrastructure Issue Reporting and Resolution Platform \\
  Project Proposal for UCS503P \\
  Submitted to: \BvrLabInstructorValue \\
  \bfseries Thapar Institute of Engineering and Technology%
}

\author{%
  Lavanya Sharma, COE%
    \thanks{Roll No: \texttt{[ROLL NUMBER]}, Email: \texttt{[EMAIL]}}
  \and
  Parismita Thapa, COE%
    \thanks{Roll No: \texttt{[1024030309]}, Email: \texttt{pthapa\_be24@thapar.edu}}
  \and
  Rudra Yadav, COE%
    \thanks{Roll No: \texttt{[1024030389]}, Email: \texttt{ryadav\_be24@thapar.edu}}
  \and
  Nadeem Esrar, COE%
    \thanks{Roll No: \texttt{[ROLL NUMBER]}, Email: \texttt{[EMAIL]}}
}

\date{\ August 18, 2026}

\begin{document}
\maketitle


% =======================================================
\section{Higher-Order Goal
  \BvrHint{\ldots secondary framing}}
% =======================================================
This project applies standard software engineering practice, including
requirement analysis, system design, and structured development, to a real
campus problem. The broader motivation is to improve campus infrastructure
governance by reducing the time and effort required to resolve civic
issues reported by students and faculty. The immediate engineering
objective is to design and implement a measurable, deployable software
solution for efficient, department-routed issue resolution across the
Thapar campus, including academic blocks, hostels, the library,
cafeterias, the sports complex, and other common areas.


% =======================================================
\section{Time-to-value
  \BvrHint{\ldots primary heuristic}}
% =======================================================
\begin{itemize}[leftmargin=0.7cm]
  \item We will deliver meaningful increments quickly by prototyping the
    core reporting-and-tracking workflow and validating it with a small
    group of students and hostel residents within a short iteration
    window.
  \item We will prioritize fast feedback over perfection by using weekly
    iterations, deferring features like auto-escalation and notifications
    until the core flow is stable, and targeting CI-backed automation
    early.
  \item Success is defined by what can be delivered and measured in the
    first iteration, a working submission-to-status flow, then improved
    based on user feedback before adding duplicate detection and
    escalation.
\end{itemize}


% =======================================================
\section{Problem Statement
  \BvrHint{\ldots the campus context}}
% =======================================================
Campus infrastructure faults -- broken equipment, plumbing leaks, faulty
wiring, damaged furniture, non-functional washrooms, unsafe walkways -- are
today reported informally: through hostel WhatsApp groups, verbal
complaints to wardens, or ad-hoc emails to department staff. This produces
several concrete problems:

\begin{itemize}[leftmargin=0.7cm]
  \item \textbf{No single source of truth:} the same issue is reported
    independently to multiple people, with no shared record of what has
    already been flagged or fixed.
  \item \textbf{No accountability or tracking:} a reporter has no way to
    check whether their complaint was even received, let alone what stage
    it is at.
  \item \textbf{No severity-aware prioritization:} a loose wire in a
    washroom and a broken chair in a common room compete for the same
    informal attention, with no structural way to surface safety-critical
    issues first.
  \item \textbf{Redundant effort:} without a shared, searchable list of
    open issues, students and faculty re-report problems that are already
    known, and administrators re-investigate what has already been
    logged.
  \item \textbf{No institutional visibility:} there is no aggregate view
    -- by building, department, or category -- that lets facilities
    management identify recurring or high-impact problem areas.
\end{itemize}

\textbf{Why this matters now:} TIET's campus spans multiple hostels,
academic blocks, and shared facilities maintained by different
departments. As the reporting channel is informal and department-specific,
the effective bottleneck is not maintenance capacity but \emph{issue
visibility and routing}. A structured reporting and tracking system
directly addresses that bottleneck without requiring any change to how
maintenance work itself is carried out.


% =======================================================
\section{Objectives
  \BvrHint{\ldots what the system is for}}
% =======================================================
\begin{itemize}[leftmargin=0.7cm]
  \item Give students and faculty a single, low-friction channel to
    report infrastructure issues with a photo, category, and precise
    location.
  \item Automatically route each report to the right department, while
    still allowing a manual override when the auto-suggested category is
    wrong.
  \item Prevent duplicate reports of the same problem from fragmenting
    attention, by detecting likely duplicates at submission time and
    letting users upvote an existing report instead.
  \item Ensure high-impact or safety-critical issues are not left to
    languish in a queue by escalating priority automatically once an
    issue crosses a configurable upvote threshold.
  \item Close the feedback loop for reporters through status-change and
    escalation notifications, so that ``I reported it and heard nothing''
    is no longer the default experience.
  \item Give department administrators and a Super-Admin department-scoped
    and institution-wide visibility respectively, so ownership of an
    issue is always clear.
\end{itemize}


% =======================================================
\section{Proposed Solution
  \BvrHint{\ldots system overview}}
% =======================================================

\subsection{User Roles}
Access is role-based, with two reporting roles and two administrative
roles:

\begin{center}
\begin{tabular}{@{}lccccl@{}}
  \toprule
  \textbf{Role} & \textbf{Report} & \textbf{View all}
    & \textbf{Upvote} & \textbf{Change status} & \textbf{Scope} \\
  \midrule
  Student        & \checkmark & \checkmark & \checkmark & --         & -- \\
  Faculty        & \checkmark & \checkmark & \checkmark & --         & -- \\
  Administrator  & --         & \checkmark & --         & \checkmark & Own department \\
  Super-Admin    & --         & \checkmark & --         & \checkmark & All departments \\
  \bottomrule
\end{tabular}
\end{center}

Students and faculty are treated identically: both can submit reports,
browse the full (read-only) list of reported issues, and upvote an
existing report. Administrators are scoped to a single department/category
and move issues through their lifecycle. The Super-Admin has full
cross-department visibility and control, and is additionally the only
role notified on every escalation event regardless of which department
owns the issue -- this gives an institutional oversight layer even when
an owning Administrator has not yet acted.

\subsection{Core Features}
\begin{itemize}[leftmargin=0.7cm]
  \item \textbf{Structured reporting:} photo upload, category selection
    (with automatic department suggestion, manually overridable by the
    reporter or an admin), and location entry via a
    \emph{Building $\rightarrow$ Floor $\rightarrow$ Area/Room} hierarchy
    (dropdowns plus free text), rather than GPS or a map pin.
  \item \textbf{Duplicate detection at submission time:} the system
    checks the new report against \emph{open} issues in the same
    location/department bucket, using location match, category match, and
    NLP-based description similarity. If a likely duplicate is found, the
    reporter is shown it and prompted to upvote instead of filing a new
    report.
  \item \textbf{Auto-escalation:} when an open issue's upvote count
    crosses a configurable, per-category threshold, it is flagged
    \emph{High Priority}, surfaced at the top of the owning
    Administrator's and the Super-Admin's dashboards, and triggers a
    Super-Admin notification. Escalation is purely a priority/visibility
    flag; it does not change the underlying status lifecycle.
  \item \textbf{Status tracking:} every issue moves through
    \texttt{Reported $\rightarrow$ Ongoing $\rightarrow$ Finished},
    updated only by the Administrator of the owning department or the
    Super-Admin.
  \item \textbf{Notifications:} in-app notifications for status changes
    (to the original reporter), upvote/duplicate activity (to the
    original reporter, so they know their issue is gaining traction), new
    reports landing in a department (to that department's Administrator),
    and escalations (to the Super-Admin).
\end{itemize}

\subsection{Core Workflow}
\begin{enumerate}[leftmargin=0.7cm]
  \item A student or faculty member submits an issue with a photo,
    category, and structured location.
  \item The system auto-suggests a department and simultaneously checks
    for likely duplicates among open issues in that location/department
    bucket.
  \item If a likely duplicate exists, the reporter is shown it and can
    upvote it instead of creating a new record; if not, a new issue is
    created in \texttt{Reported} status and routed to the relevant
    Administrator, who is notified.
  \item As upvotes accumulate, the system checks the issue's count
    against its category's escalation threshold on every upvote event;
    on crossing it, the issue is marked \emph{High Priority} and the
    Super-Admin is notified.
  \item The owning Administrator (or Super-Admin) moves the issue through
    \texttt{Reported $\rightarrow$ Ongoing $\rightarrow$ Finished}; the
    original reporter is notified at each transition.
\end{enumerate}

\subsection{Operational Constraints for an Educational, Tier-2/3 Setting}
\begin{itemize}[leftmargin=0.7cm]
  \item Keep the solution usable on low-to-mid range student devices and
    typical hostel/campus network conditions, via responsive web design
    rather than a heavier native app.
  \item Avoid dependence on research-grade algorithms for duplicate
    detection or routing; prioritize workflow correctness, routing
    accuracy, and system reliability over algorithmic novelty.
  \item Design for deployability using common managed hosting and a
    managed database, since the team does not have access to dedicated
    operations infrastructure within the course timeline.
\end{itemize}


% =======================================================
\section{Solution Approach
  \BvrHint{\ldots engineering focus}}
% =======================================================
This is an engineering course project, so the emphasis is on scalable
workflow and system-engineering concerns, duplicate handling, service
architecture, and delivery pipeline, rather than on novel algorithm
research.

\subsection{Duplicate/Quality Assistance}
Duplicate detection is heuristic-based, combining three signals so that
no single weak signal can produce a false match on its own:

\begin{enumerate}[leftmargin=0.7cm]
  \item \textbf{Location match} -- exact match on building and floor,
    fuzzy match on the free-text area/room field.
  \item \textbf{Department/category match} -- the new report is only
    compared against issues already routed to the same department.
  \item \textbf{Description similarity} -- a simple similarity measure
    (e.g., token overlap or TF-IDF cosine similarity with a threshold,
    without any claim to state-of-the-art accuracy) between the new
    description and existing open descriptions within the matched
    location/department bucket.
\end{enumerate}

The comparison set is restricted to issues still open (not
\texttt{Finished}), on the reasoning that a resolved issue recurring is a
\emph{new} problem worth a fresh report, not a duplicate. The system
surfaces the top candidate to the reporter for a human decision -- it
never fully auto-merges, auto-deletes, or auto-rejects a submission --
keeping a human in the loop for a step that directly affects data
quality.

\subsection{Web Architecture}
A typical 3-tier web architecture:
\begin{itemize}[leftmargin=0.7cm]
  \item \textbf{Frontend:} React, a responsive UI for both reporters
    (submission form, status view) and administrators (dashboard),
    usable on low-to-mid range student devices.
  \item \textbf{Backend API:} Node.js/Express (or Flask as a fallback),
    exposing a REST API for authentication, submission, routing, status
    transitions, and escalation logic.
  \item \textbf{Database:} PostgreSQL, chosen for relational integrity
    between issues, departments, users, upvotes, and notifications, and
    for straightforward indexing on the fields duplicate detection and
    dashboards query most (location, department, status).
\end{itemize}
Supporting infrastructure: Cloudinary or AWS S3 for image attachments,
kept out of the primary database; JWT-based sessions with role-based
access control enforced at the API layer for all four roles; and
notification delivery via polling rather than WebSockets, to keep the
MVP's infrastructure and failure modes simple.

\subsection{Data Model Highlights}
\begin{itemize}[leftmargin=0.7cm]
  \item An \texttt{issues} table carries category, department,
    structured location fields, status, an upvote count, and a
    \texttt{priority} field (\texttt{Normal}/\texttt{High}) set by the
    escalation check.
  \item An \texttt{escalation\_threshold} value, configurable by the
    Super-Admin and tunable per category, is checked against the upvote
    count on every upvote event.
  \item A \texttt{notifications} table (\texttt{user\_id},
    \texttt{issue\_id}, \texttt{type}, \texttt{message},
    \texttt{read\_status}, \texttt{timestamp}) backs all four
    notification types described above.
  \item Duplicate-location matching is a structured field comparison
    (exact building/floor, fuzzy area/room text) rather than a
    geospatial radius calculation -- a deliberate trade-off, since GPS
    is unreliable indoors and map APIs add cost and complexity the MVP
    does not need.
\end{itemize}

\subsection{CI/CD and Fast Delivery Enabler}
CI/CD will be used to:
\begin{itemize}[leftmargin=0.7cm]
  \item Automatically build and run tests (backend unit tests, linting)
    on every change.
  \item Produce repeatable, versioned deployment artifacts to a staging
    environment.
  \item Enable safe, frequent releases of incremental features, in line
    with the weekly iteration cadence described under Time-to-value.
\end{itemize}


% =======================================================
\section{Auto-Escalation Design
  \BvrHint{\ldots priority-aware triage}}
\label{sec:auto-escalation}
% =======================================================
A single flat upvote threshold across every category would either
(a) be too high for safety-critical issues, which should not have to
wait for popularity to earn attention, or (b) be too low for
high-traffic areas, where ``High Priority'' would stop meaning anything.
The design instead makes the threshold configurable per category, set by
the Super-Admin:

\begin{itemize}[leftmargin=0.7cm]
  \item \textbf{Safety-critical} (electrical, structural, fire safety):
    a low threshold ($\sim$2--3 upvotes) -- severity, not volume, should
    drive escalation here.
  \item \textbf{High-traffic common areas} (library, cafeteria, sports
    complex): a higher threshold ($\sim$8--12 upvotes), since large
    natural visibility means genuine problems accumulate upvotes
    quickly on their own.
  \item \textbf{Hostel/block-specific issues:} an intermediate threshold
    ($\sim$5--8 upvotes), reflecting a smaller audience per location.
  \item \textbf{Low-severity/cosmetic issues} (paint, furniture): a
    higher bar ($\sim$15+ upvotes), so that ``High Priority'' is not
    diluted.
\end{itemize}

For the MVP, a single \textbf{default threshold of 5 upvotes} is used
across all categories, since no real usage data from TIET exists yet to
calibrate per-category values responsibly. This default is treated as a
known, temporary simplification, explicitly intended to be reviewed and
tuned per category after the pilot, once real reporting-volume data is
available -- not as a final design decision.


% =======================================================
\section{Evaluation Criteria
  \BvrHint{\ldots measurable and attributable}}
% =======================================================
The system's original README does not specify success metrics; the
following are proposed based on the features actually being built, so
that each stated feature has a corresponding measurable outcome.

\subsection{Primary Metrics}
\label{sec:eval-primary}
\begin{itemize}[leftmargin=0.7cm]
  \item \textbf{Time-to-first-action (TTFA):} time between an issue
    entering \texttt{Reported} and its first transition to
    \texttt{Ongoing}, measured from database timestamps. Target: median
    TTFA within a defined pilot-window bound (to be set from pilot
    baseline data, since no institutional baseline currently exists).
  \item \textbf{Duplicate-detection precision:} of the duplicate
    candidates surfaced to reporters, the proportion that reporters
    agree are genuine duplicates (accepted upvote vs.\ rejected and
    filed as new). This is directly measurable from user choices at
    submission time and avoids needing a hand-labelled test set.
\end{itemize}

\subsection{Secondary Metrics}
\begin{itemize}[leftmargin=0.7cm]
  \item \textbf{Escalation precision:} proportion of auto-escalated
    issues that Administrators, on review, agree warranted High
    Priority -- a check on whether the MVP's flat threshold is behaving
    reasonably ahead of per-category tuning.
  \item \textbf{Resolution time by category:} time from
    \texttt{Reported} to \texttt{Finished}, broken down by category, to
    identify systemic bottlenecks by department rather than only
    reporting an institution-wide average.
  \item \textbf{Duplicate deflection rate:} proportion of submissions
    that end as an upvote on an existing issue rather than a new report
    -- a direct measure of how much redundant reporting is being
    prevented.
  \item \textbf{Notification effectiveness:} proportion of
    status-change notifications that are read (via
    \texttt{read\_status}) within a defined window, as a proxy for
    whether the feedback loop is actually closing the trust gap
    identified in the problem statement.
\end{itemize}

\subsection{Pilot Validation Plan}
\label{sec:pilot-validation}
\begin{itemize}[leftmargin=0.7cm]
  \item Run the pilot within a bounded scope -- e.g., one or two hostel
    blocks plus one academic block -- rather than campus-wide, to keep
    the reporting volume interpretable within the course timeline.
  \item Recruit a small set of student/faculty reporters and, for
    Administrator accounts, either real departmental staff (if
    available) or team-simulated department roles.
  \item Collect an informal baseline (e.g., via a short survey on how
    long informally-reported issues currently take to get attention) to
    give the primary/secondary metrics a comparison point, since no
    existing institutional baseline is being tracked today.
\end{itemize}


% =======================================================
\section{Scalability and Deployment
  \BvrHint{\ldots beyond the pilot}}
% =======================================================
\begin{enumerate}[leftmargin=0.7cm]
  \item \textbf{Theoretical scalability:}
    \begin{itemize}[leftmargin=0.7cm]
      \item Decoupled frontend/backend, communicating over a stateless
        REST API, so either layer can scale independently.
      \item Indexing on the fields most queried by duplicate detection
        and dashboards (location, department, status), plus pagination
        on list views, so query cost does not grow linearly with total
        issue volume as adoption increases beyond the pilot.
      \item Escalation and duplicate-detection checks scoped to a
        bounded bucket (same department/location, open issues only)
        rather than the full issue table, so their cost does not grow
        with total historical volume either.
    \end{itemize}
  \item \textbf{Practical deployability:}
    \begin{itemize}[leftmargin=0.7cm]
      \item Common managed hosting for the backend and a managed
        PostgreSQL instance, avoiding bespoke infrastructure.
      \item Object storage (Cloudinary/S3) for images from day one, so
        image volume growth does not pressure the primary database.
      \item CI/CD pipeline for staging/production, as detailed under
        CI/CD and Fast Delivery Enabler.
    \end{itemize}
\end{enumerate}


% =======================================================
\section{Engine Availability Heuristic}
% =======================================================
A mature set of ``engines'' (tooling/services) is available to
implement this as a web-based project, so the team is not building
infrastructure from scratch:

\begin{itemize}[leftmargin=0.7cm]
  \item Web framework for REST APIs (Express/Flask) and a responsive
    frontend framework (React) for templated/reactive pages.
  \item A managed relational database (PostgreSQL).
  \item Token-based authentication (JWT) with role-based access
    control.
  \item Object storage for photo attachments (Cloudinary/AWS S3, with
    local storage as a pilot fallback).
  \item A test framework for backend unit/integration tests.
  \item A CI runner (e.g., GitHub Actions) and a staging deployment
    target.
\end{itemize}

We will use these mature, off-the-shelf components to focus engineering
effort on workflow design, duplicate/escalation logic, and measurement,
rather than on inventing new infrastructure.


% =======================================================
\section{Project Scope and Deliverables
  \BvrHint{\ldots phased delivery}}
% =======================================================

\subsection{MVP Deliverables}
\begin{itemize}[leftmargin=0.7cm]
  \item Citizen-side (student/faculty) issue submission form: photo,
    category with auto-suggested department, structured location entry.
  \item Duplicate-candidate check at submission time (location +
    department + description similarity), with an upvote path.
  \item Department-scoped Administrator dashboard: view assigned
    issues, update status through \texttt{Reported $\rightarrow$
    Ongoing $\rightarrow$ Finished}.
  \item Super-Admin dashboard: cross-department visibility,
    high-priority issues surfaced at the top.
  \item Auto-escalation with the flat MVP default threshold
    (5 upvotes).
  \item In-app notifications for all four event types (status change,
    upvote/duplicate, new-report-to-admin, escalation-to-Super-Admin),
    via polling.
  \item JWT-based auth with role-based access control for all four
    roles.
  \item CI: build, backend unit tests, linting on every change; CD to a
    staging environment.
\end{itemize}

\subsection{Post-MVP / Subsequent Deliverables}
\begin{itemize}[leftmargin=0.7cm]
  \item Per-category configurable escalation thresholds (replacing the
    flat MVP default), calibrated using real pilot upvote-volume data.
  \item Improved duplicate-detection similarity (e.g., embedding-based
    similarity instead of pure token overlap/TF-IDF) once enough real
    description data exists to evaluate it meaningfully.
  \item Aggregate analytics for facilities management: recurring-issue
    hotspots by building/category, resolution-time trends over time.
  \item Optional email/SMS notification channel alongside in-app
    notifications.
  \item Search/filtering improvements on the public issue list as
    report volume grows.
\end{itemize}


% =======================================================
\section{Risks and Mitigations
  \BvrHint{\ldots anticipated challenges}}
% =======================================================
\begin{itemize}[leftmargin=0.7cm]
  \item \textbf{Duplicate false positives/negatives:} an overly
    aggressive similarity threshold could push genuinely distinct
    issues together, discouraging legitimate reports; an overly
    conservative one would fail to deflect true duplicates.
    \emph{Mitigation:} always show the candidate for a human decision
    rather than auto-merging, and treat the duplicate-detection
    precision metric (Section~\ref{sec:eval-primary}) as a live signal
    for tuning the threshold.
  \item \textbf{Upvote gaming:} without a one-vote-per-user-per-issue
    constraint, upvote counts (and therefore escalation) could be
    manipulated. \emph{Mitigation:} enforce a unique
    (\texttt{user\_id}, \texttt{issue\_id}) constraint at the database
    level on upvotes.
  \item \textbf{MVP threshold miscalibration:} the flat default-5
    threshold is known to be a simplification and may over- or
    under-escalate relative to a category's real severity/volume
    profile. \emph{Mitigation:} explicitly scoped as an MVP limitation
    (Section~\ref{sec:auto-escalation}), with per-category tuning as a
    defined post-pilot deliverable rather than an afterthought.
  \item \textbf{Low pilot adoption:} a reporting tool is only useful if
    people use it instead of the informal channels it is meant to
    replace. \emph{Mitigation:} keep the submission flow to a minimal
    number of fields, and scope the pilot narrowly
    (Section~\ref{sec:pilot-validation}) so the team can actively drive
    adoption within a small, trackable population.
  \item \textbf{Administrator adoption:} if department staff do not
    update statuses, the notification loop and TTFA metric both break
    down. \emph{Mitigation:} keep the Administrator update action to a
    single dropdown change, and simulate Administrator roles within the
    team if real departmental staff are unavailable during the pilot
    window.
  \item \textbf{Image storage cost/volume:} unrestricted photo uploads
    could grow storage cost quickly. \emph{Mitigation:} client-side
    compression before upload and a per-report image-count/size cap.
\end{itemize}


% =======================================================
\section{Summary}
% =======================================================
This project proposes Civic Issue Reporter, a role-based web platform
for reporting, routing, and resolving campus infrastructure issues at
TIET. Its core engineering contributions are a submission flow that
actively suppresses duplicate reports rather than merely listing them,
and a priority-escalation mechanism designed from the outset to be
severity-aware rather than purely volume-driven, even though the MVP
ships with a single calibratable default. The proposal defines primary
and secondary metrics directly tied to these mechanisms
(time-to-first-action, duplicate-detection precision, escalation
precision) so that the pilot produces the data needed to move from the
MVP's flat defaults to the tuned, per-category design that motivated the
system in the first place.

\end{document}
